\documentclass[10pt,a4paper]{amsart}
\usepackage[margin=19mm]{geometry}
\usepackage{amsmath,amssymb,mathtools}
\usepackage[scr=rsfs,cal=euler]{mathalfa}
\usepackage{xfrac}
\usepackage[unicode,hidelinks,bookmarks=false]{hyperref}
\newcommand{\ie}{i.e.\ }
\newcommand{\cf}{cf.\ }
\newcommand{\resp}{respectively\ }
\newcommand{\Subsection}{\subsection}
\providecommand{\nobreakdash}{\nobreak\hbox{-}}
\setlength{\emergencystretch}{2em}
\multlinegap0pt
\usepackage{amssymb}
\usepackage[shortlabels]{enumitem}
\usepackage{xcolor}
\usepackage{mathtools}
\newtheorem{prop}{Proposition}
\newcommand{\Hom}{\operatorname{Hom}}
\renewcommand{\Im}{\operatorname{Im}}
\newcommand{\Ker}{\operatorname{Ker}}
\newcommand{\ML}{\operatorname{ML}}
\newcommand{\Mod}{\operatorname{Mod}}
\newcommand{\N}{\mathbb{N}}
\newcommand{\Pro}{\operatorname{Pro}}
\newcommand{\TopMod}{\operatorname{TopMod}}
\renewcommand{\bar}{\overline}
\newcommand{\colim}{\operatornamewithlimits{\rm{colim}}}
\newcommand{\eq}[2]{\begin{equation}\label{#1}#2 \end{equation}}
\renewcommand{\lim}{\operatornamewithlimits{\rm{lim}}}
\renewcommand{\phi}{{\varphi}}
\newcommand{\ra}{\rightarrow} % Fei added
\newcommand{\xra}{\xrightarrow}
\title{Notes on Coherent six-functor formalisms: Pro vs Solid}
\author{Fei Ren}
\date{Source: arXiv:2506.21082v3 (CC BY 4.0)}
\begin{document}
\maketitle

\noindent\emph{Source and licence.} Notes on Coherent six-functor formalisms: Pro vs Solid. Version arXiv:2506.21082v3 (CC BY 4.0), distributed by arXiv under the Creative Commons Attribution 4.0 International licence, http://creativecommons.org/licenses/by/4.0/, as recorded in the arXiv licence metadata for this exact version and shown on the abstract page https://arxiv.org/abs/2506.21082v3. Full licence notice: \texttt{licenses/arxiv-2506.21082-CC-BY-4.0.txt}.\par\smallskip

\noindent\textit{Editorial scope.} Counted unit: the article's prop Phi1FF (source0.tex lines 545--550). The article's complete original proof of it is delivered with it (source0.tex lines 552--598, one \texttt{proof} environment of 47 lines). No mathematics is written, repaired or re-derived: the statement and the proof are the article's own lines, in the article's own notation, and no retained line is substituted or re-flowed.  No further source-local context is needed: every cross-reference of every retained line resolves inside the two delivered spans.  Nothing was omitted from any retained span, so the package carries no omission record.  No retained line contains a citation, so the wrapper carries no bibliography and no bibliography entry.  No retained line contains a figure environment, an image-inclusion command or drawing code, so no image is delivered and the package declares no asset dependency.  Editorial formatting only, in addition: an AMS \texttt{amsart} preamble that restates the article's own declarations and macros used by the retained lines, namely the environments \texttt{prop} on separate counters, together with the standard packages those lines need.  The retained environments print in this document as Proposition 1 in order of appearance, whereas the article prints the same items with the numbering of its own full document.  The complete native source of the article as distributed in the arXiv e-print for this exact version is bundled unchanged as \texttt{sources/180527/source0.tex}.
\begin{prop}\label{Phi1FF}
The following restriction of $\Phi_1$
\eq{{Phi1ML}}{
\Pro^{\ML}_{\N}\Mod(A)\subset \Pro_\N\Mod(A)\xra{\Phi_1} \TopMod(A),\quad``\lim " M_i\ra \lim M_i}
is fully faithful.
\end{prop}

\begin{proof}
We prove something slightly stronger:
Let $``\lim_{i\in \N}" M_i\in \Pro^{\ML}_{\N}\Mod(A), ``\lim_J" N_j\in \Pro\Mod(A)$, the category of pro-$A$-modules with any index set. Then
\begin{align*} 
\Hom_{\Pro\Mod(A)}(``\lim_{i\in \N}" M_i,``\lim_J"N_j)
&\stackrel{(1)}{=}\lim_J \colim_{i\in \N} \Hom_{\Mod(A)}(M_i,N_j)
\\
&\stackrel{(2)}{=}\lim_J \Hom_{\TopMod(A)}(\lim_{i\in \N} M_i,N_j)\\
&\stackrel{(3)}{=}\Hom_{\TopMod(A)}(\lim_{i\in \N} M_i,\lim_J N_j).
\end{align*}
Here 
(1) follows from the definition of the pro-category.     

(2) 
We show that the canonical map 
$$\colim_{i\in \N} \Hom_{\Mod(A)}(M_i, N_j) \to \Hom_{\TopMod(A)}(\lim_{i\in \N} M_i, N_j)$$
is a bijection.

\textit{Surjectivity}.
Let $f: \lim_{i\in \N} M_i \to N_j$ be a continuous map. Since the target space $N_j$ is discrete, $\{0\}$ is open, so $\Ker(f)$ is an open neighborhood of $0$ in $\lim_{i\in \N} M_i$. 
By the definition of the inverse limit topology, a basic open neighborhood of $0$ is a finite intersection of sets of the form $\Ker(p_i)$. Thus, there exists a finite subset $I' \subset \N$ such that $\bigcap_{i \in I'} \Ker(p_i) \subset \Ker(f)$. Because $\N$ is totally ordered, $I'$ has a maximum element $k$. For any $i \in I'$, $p_i$ factors through $p_k$, meaning $\Ker(p_k) \subset \Ker(p_i)$. Thus $\Ker(p_k) \subset \Ker(f)$.
This inclusion means $f$ factors through the image of $p_k$, yielding a well-defined map 
$$\bar{f}: \Im(p_k) \to N_j.$$
Next, we use the Mittag-Leffler condition. Because $``\lim_{i\in \N}" M_i$ is Mittag-Leffler, there exists an index $l \ge k$ such that for all $m \ge l$, $\Im(\phi_{mk}) = \Im(\phi_{lk})$. 
The Mittag-Leffler condition, together with the countable index set, guarantees that the 
projection map $p_k:\lim_{i\in \N} M_i\ra M_k$ maps
exactly onto this stable image, namely 
$$\Im(p_k) = \Im(\phi_{lk}).$$
Therefore, our map $\bar{f}$ can be rewritten as
$$\bar f: \Im(\phi_{lk}) \ra N_j.$$
Define
$$f_l = \bar{f} \circ \phi_{lk}: M_l \to N_j.$$
This $f_l\in \colim_{i\in \N} \Hom_{\Mod(A)}(M_i, N_j)$ clearly maps to $f$ under the canonical map, proving surjectivity.


\textit{Injectivity.}
Suppose two maps $g_1, g_2: M_k \to N_j$ have the same image in $\Hom_{\TopMod(A)}(\lim_{i\in \N} M_i, N_j)$. Let $g = g_1 - g_2$. Then 
$$g \circ p_k = 0.$$
By the same Mittag-Leffler stabilization argument, there exists an $l \ge k$ such that $\Im(\phi_{lk}) = \Im(p_k)$. 
Therefore, $g \circ \phi_{lk} = 0$, which means 
$$g_1 \circ \phi_{lk} = g_2 \circ \phi_{lk}.$$
This implies $g_1$ and $g_2$ define the same element in the colimit $\colim_{i\in \N} \Hom_{\Mod(A)}(M_i, N_j)$, proving injectivity.

(3)
follows from the universal property of $\lim_J$.

\end{proof}

\end{document}
